\section{Microcontrolador}

O ESP32 \cite{esp32_datasheet} é um microcontrolador de duplo núcleo (Xtensa LX6, 240 MHz) desenvolvido pela Espressif Systems, dotado de conectividade Wi-Fi 802.11 b/g/n e Bluetooth 4.2/BLE integrados, ADC de 12 bits com 18 canais, memória flash de 4 MB e suporte a protocolo SPI, I2C e UART. Sua ampla disponibilidade, baixo custo e capacidade de processamento em tempo real tornam adequado tanto para a aquisição e transmissão dos sinais sEMG (ESP32 transmissor) quanto para o processamento e controle da prótese (ESP32 receptor).
A comunicação entre os dois módulos ESP32 pode ser realizada via protocolo ESP-NOW, que oferece baixa latência, baixo consumo energético e conexão ponto a ponto sem necessidade de infraestrutura de rede, sendo ideal para aplicações em tempo real como o controle protético

